% Compile twice: lualatex 20001192.tex
\documentclass[11pt]{article}
\usepackage[margin=25mm]{geometry}
\usepackage{fontspec,unicode-math}
\setmainfont{Latin Modern Roman}
\setmathfont{Latin Modern Math}
\usepackage{amsmath,amssymb,mathtools}
\usepackage{xurl,hyperref}
\hypersetup{colorlinks=true,urlcolor=blue,linkcolor=blue}
\setcounter{secnumdepth}{0}
\setlength{\parindent}{0pt}
\setlength{\parskip}{5pt}
\setlength{\emergencystretch}{3em}
\newcommand{\sref}[2]{\hyperlink{source:#1:#2}{[S#2]}}
\newcommand{\cataloguescope}[1]{\paragraph{Source question.}#1}
\newcommand{\floor}[1]{\left\lfloor#1\right\rfloor}
\begin{document}
% UnsolvedMath ID 20001192; source catalogue code AIM-COMPUTATION-0030
\setcounter{section}{20001191}
\section{Holonomic Bayesian normalizers with certified posterior moments}\label{20001192}
{\small\sffamily UnsolvedMath ID 20001192 · Computer Science}
\subsection{Definitions and mathematical statement}

\paragraph{Shared notation.}Standard number systems are denoted by $\mathbb N,\mathbb Z,\mathbb Q,\mathbb R,\mathbb C$; index ranges, computational models and specialized notation follow the cited problem source and the definitions below.


For an unnormalized density $g(\theta)$, its normalizing constant is $Z=\int g(\theta)\,d\theta$ over the specified domain, when this integral exists. Bayes' formula normalizes the product of a prior density and likelihood. The source asks for algebraic methods for exact expressions and approximations of the resulting integrals.
\paragraph{Statement recorded in the supplied source.}
Question 16. Computing normalization constants and posterior distributions for continuous random variables involves the computation of integrals involving rational and transcendental funtions. Can algebraic techniques be used to give exact expressions for these integrals or help give approximations for them? 7 \sref{20001192}{2}
\paragraph{Scope qualification.}
No universal exact symbolic form or particular approximation guarantee is stipulated in the source.
\paragraph{Source status.}
The supplied catalogue labels this record partially solved. Its source-status metadata was checked against the supplied snapshot on 2026-10-11. This import does not independently establish current open status or mathematical correctness.
\subsection{Short English statement}
Can algebra help compute integrals needed to normalize continuous probability models and their posterior distributions?
\subsection{Sources}
\begin{itemize}
\item[\hypertarget{source:20001192:1}{S1}] Ulam AI and the UnsolvedMath Contributors, UnsolvedMath: A Curated Collection of Open Mathematics Problems, record 20001192 (AIM-COMPUTATION-0030). Catalogue curation and metadata: CC BY 4.0. The supplied source question is retained; definitions and the English summary are editorial additions.
\url{https://huggingface.co/datasets/ulamai/UnsolvedMath}.
\item[\hypertarget{source:20001192:2}{S2}] AIM Workshop Problem Lists, original problem source linked by catalogue record AIM-COMPUTATION-0030
\url{https://aimath.org/WWN/compalgstat/compalgstat.pdf}.
\end{itemize}
\paragraph{Preparation and review.}
Original question attributed to the cited source and UnsolvedMath contributors. Codex (GPT-6; exact serving revision not exposed) assisted with definitions, English summaries, source inspection and formatting. No new human mathematical authorship or independent proof verification is asserted. Canonical ID 20001192 and source code AIM-COMPUTATION-0030 checked against the complete supplied registry and existing repository statements; known flagged repetitions and obvious identical conjectures excluded. Full mathematical-equivalence/literature review has not been performed. Statement preparation and final review followed contribution prompts C and D; automated checks verify structure and identity, not open status.
\label{end:20001192}
\end{document}
